\usepackage{amsmath,amssymb}
\usepackage{xcolor}
\usepackage{tikz}
\usetikzlibrary{arrows.meta,backgrounds,positioning}

\definecolor{inputfill}{HTML}{BFE3C0}
\definecolor{inputdraw}{HTML}{24592B}
\definecolor{outputfill}{HTML}{F1BBB5}
\definecolor{outputdraw}{HTML}{72231F}
\definecolor{edgecolor}{HTML}{2B4472}
\definecolor{textcolor}{HTML}{172033}

\tikzset{
  >=Latex,
  input neuron/.style={
    circle,
    draw=inputdraw,
    fill=inputfill,
    line width=0.85pt,
    minimum size=10mm,
    inner sep=0pt,
    font=\bfseries\boldmath\normalsize
  },
  hidden neuron/.style={
    circle,
    draw=edgecolor,
    fill=hiddenfill,
    line width=0.9pt,
    minimum size=10mm,
    inner sep=0pt,
    font=\bfseries\boldmath\normalsize
  },
  output neuron/.style={
    circle,
    draw=outputdraw,
    fill=outputfill,
    line width=0.85pt,
    minimum size=10mm,
    inner sep=0pt,
    font=\bfseries\boldmath\normalsize
  },
  connection/.style={
    draw=edgecolor!48,
    line width=0.50pt,
    -{Latex[length=1.8mm,width=1.25mm]}
  },
  layer title/.style={
    font=\bfseries\boldmath\small,
    text=textcolor,
    align=center
  },
  parameter/.style={
    font=\bfseries\small,
    text=textcolor,
    align=center
  }
}

\definecolor{hiddenfill}{HTML}{CBC9F2}

\newcommand{\connectinputhidden}[2]{%
  \begin{scope}[on background layer]
    \foreach \i in {#1}{
      \foreach \h in {#2}{
        \draw[connection] (I-\i) -- (H-\h);
      }
    }
  \end{scope}
}

\newcommand{\connecthiddenoutput}[2]{%
  \begin{scope}[on background layer]
    \foreach \h in {#1}{
      \foreach \j in {#2}{
        \draw[connection] (H-\h) -- (O-\j);
      }
    }
  \end{scope}
}
